% Generated draft: requires mathematical and source review.
\documentclass[11pt]{article}
\usepackage[margin=25mm]{geometry}
\usepackage{fontspec}
\setmainfont{Latin Modern Roman}
\usepackage{amsmath,amssymb,mathtools,unicode-math}
\setmathfont{Latin Modern Math}
\usepackage{enumitem,needspace,xurl,hyperref}
\setlength{\parindent}{0pt}
\setlength{\parskip}{5pt}
\setlength{\emergencystretch}{3em}
\newcommand{\N}{\mathbb{N}}
\newcommand{\Z}{\mathbb{Z}}
\newcommand{\Q}{\mathbb{Q}}
\newcommand{\R}{\mathbb{R}}
\newcommand{\C}{\mathbb{C}}
\newcommand{\F}{\mathbb{F}}
\newcommand{\A}{\mathbb{A}}
\newcommand{\PP}{\mathbb{P}}
\newcommand{\E}{\mathbb{E}}
\newcommand{\eps}{\varepsilon}
\newcommand{\dd}{\,\mathrm{d}}
\begin{document}
\section*{UnsolvedMath ID: 30001122}
\textbf{Source problem code:} catalogue entry
\subsection{Definitions and mathematical statement}
An a posteriori error bound is an upper bound on an approximation error obtained from the computed approximation and the prescribed problem data. An associated adaptive mesh algorithm uses error information to select changes to the computational mesh.

The supplied research problem is to develop reliable a posteriori error control and associated adaptive mesh algorithms for numerical approximations of the Navier--Stokes equations. The supplied literature assessment places this objective in the setting of the full nonstationary problem; the statement itself does not specify a particular Navier--Stokes system or initial-boundary value problem.

A schematic norm-bound interpretation is the following. Once a continuous problem, a solution concept, an approximation method, and an error space \(X\) have been specified, let \(U\) denote the exact solution quantities being approximated and let \(U_h\) denote their numerical approximation on a computational mesh indexed by \(h\). Seek a computable a posteriori bound \(B_h\), obtained from \(U_h\), the mesh, and the prescribed data, such that
\[
\lVert U-U_h\rVert_X \leq B_h,
\]
and develop a mesh adaptation procedure informed by that error control. This display explains a possible meaning of reliability; it is not a fully specified theorem or a claim that such a bound exists. The source does not determine the admissible problem class, the norm, the dependence of the bound on problem parameters, or the guarantees required of the adaptive procedure.

In particular, the spatial dimension and domain, compressible or incompressible formulation, coefficients, forcing, initial and boundary conditions, time interval, solution regularity and uniqueness assumptions, approximated variables, discretization, and treatment of temporal error are unspecified. Consequently, a precise universally quantified existence assertion cannot be recovered from this record without additional hypotheses. No particular choices of these missing ingredients are imposed here.
\subsection{Short English statement}
Develop computable error bounds and mesh adaptation methods for numerical solutions of the Navier--Stokes equations. The supplied assessment concerns the full nonstationary problem, but the record leaves the equations, admissible data, solution class, error measure, and algorithmic guarantees unspecified. This is therefore a research objective requiring further formulation and human mathematical review.
\subsection{Sources}
\begin{itemize}
\item Local UnsolvedMath registry, supplied complete record for permanent ID \(30001122\), titled ``Adaptive Error Control for Navier--Stokes Equations,'' in the Oberwolfach Reports Open Problems collection. Exact source problem code: \texttt{catalogue entry}. No individual proposer is supplied; the recorded proposed year is 2008.
\item Supplied report citation: \emph{Hyperbolic Conservation Laws (2009)}, Oberwolfach Reports, pp.~3139--3200, DOI \texttt{10.4171/owr/2008/56}. Original OWR report: \url{https://doi.org/10.4171/owr/2008/56}. The background describes Oberwolfach Reports as documenting mathematical workshops and their open-problem discussions. No problem-specific page or passage is supplied.
\item Additional supplied literature source: Clay Navier--Stokes problem, \url{https://www.claymath.org/millennium-problems/navier-stokes-equation/}. The record does not specify a passage establishing its connection to the requested numerical error control.
\item Literature assessment dated exactly \texttt{2026-08-21}: the supplied curation labels the problem open and states that adaptive a posteriori error control for the full nonstationary Navier--Stokes problem remains unresolved in general, while results exist for restricted stationary, linearized, or discretized settings. This is supplied dated curation, not an independently verified literature assessment. The background explicitly requests independent verification of time-sensitive status and cited claims.
\item The registry labels the statement \texttt{corrected\_verified} and reports that its mathematical content was checked against the official Oberwolfach report, with notation, context, and question form normalized without changing the source-level claim. This verification description is supplied provenance; no independent source inspection or mathematical verification is asserted here.
\end{itemize}
\end{document}
